% Generated by tools/textbook/export_listings.js from decks/plan.js and code/l15.S: do not edit.
\section{The program of this chapter}\label{sec:lab:l15}

The listings of this chapter are parts of one program, \fname{code/l15.S} in the
repository of the course, printed in full in Listing~\ref{lst:lab:l15}. The
automatic tests of the course material build it and run it on the headless
runner of \prgqt{}, with scripted input, and check the frame;
Figure~\ref{fig:lab:l15} is the frame that the test wrote.

\begin{description}[style=nextline,leftmargin=1.5em]
\item[On the screen.] Five rows of lamps. The first is the word 0x12345678 read with lw. The other four are the bytes at addresses A, A+1, A+2, A+3 read with lbu: 0x78, 0x56, 0x34, 0x12. The lowest address holds the least significant byte: little-endian.
\item[What to press.] Nothing: the program runs by itself.
\end{description}

\begin{figure}[htb]
\centering
\includegraphics[width=0.5\textwidth]{lab_l15.png}
\caption{\fname{l15.S} on \prgqt{}.}\label{fig:lab:l15}
\end{figure}

\asmfile[caption={The complete program of the chapter, \fname{code/l15.S}.},label=lst:lab:l15]{code/l15/program.S}

To run it on \prgqt{} the commands are those of Listing~\ref{lst:lab:l15:run};
Section~\ref{sec:prg:testing} gives the same steps for QEMU and for the board.

\begin{shell}[caption={Building \fname{l15.S} and running it on \prgqt{}.},label=lst:lab:l15:run]
# 1. start PRG32-QT on this computer (it listens on port 8080)

# 2. in the PRG32 folder, with ESP-IDF active: build the cartridge
python3 -m prg32 cartridge build ~/arc12/code/l15.S --portable \
    --entry-prefix l15 --name l15 --out l15.prg32

# 3. send it to PRG32-QT and start it
python3 -m prg32 esp32c6 upload-and-run l15.prg32 \
    --url http://localhost:8080 --slot cart0

# 4. look, press, compare with the previous slide
\end{shell}
